\chapter{내용 숙지 (자기 기관 학습)}

\begin{summarybox}{한눈에}
본 기관만의 \textbf{내용 숙지 학습}을 직접 만들고 운영합니다. 본문(마크다운)을 입력하고, AI 가 체크포인트와 객관식·OX·단답·서술형 문제를 자동 생성. 페이지 단위로 본문 + 문제 한 세트가 묶여 학생에게 시험지처럼 출제됩니다.
\end{summarybox}

\kfShot{07-study-content}{내용 숙지 — 본 기관 학습 + 본사 PUBLIC 학습 목록}

\section*{화면 구조}

\begin{screenbox}{학습 콘텐츠 목록 + 신규 작성}
\noindent
\begin{tabular}{@{}p{3cm}p{2cm}p{2cm}p{2cm}p{2cm}p{2.5cm}@{}}
\toprule
\textbf{제목} & \textbf{레벨} & \textbf{공개 범위} & \textbf{문항} & \textbf{상태} & \textbf{액션} \\
\midrule
삼국유사 발췌 & 러셀1 & ORG (본 기관) & 12 & active & \inacttab{비주얼 에디터} \\
조선 후기 농서 & 비트1 & PUBLIC (본사) & 20 & active & \inacttab{미리보기 (읽기)} \\
\bottomrule
\end{tabular}

\medskip
\activetab{+ 새 학습 만들기}
\end{screenbox}

\section*{여기서 할 수 있는 일}
\begin{itemize}[leftmargin=1.2em]
  \item \textbf{새 학습 만들기} — 본 기관 소속(ORG visibility)으로 자동 생성
  \item \textbf{비주얼 에디터} — 페이지 추가, 본문 마크다운 작성, AI 체크포인트 추출, AI 문제 생성, 4유형(객관식/OX/단답/서술) 직접 편집
  \item 본사가 만든 PUBLIC 학습은 \textbf{읽기 전용 미리보기}만 가능 — 수정 불가
  \item AI 자몽 차감 — 본인 기관 자몽 지갑에서 자동 차감 (Chapter 12)
\end{itemize}

\section*{자주 쓰는 흐름}
\begin{enumerate}
  \item ``+ 새 학습 만들기'' → 제목·레벨·영역 입력 → 생성
  \item 비주얼 에디터 → 페이지 추가 → 본문 입력 (또는 PDF 업로드 → AI 마크다운 변환)
  \item ``AI 체크포인트 추출'' → 출제 포인트 자동 생성 (자몽 1개)
  \item ``AI 문제 생성'' → 5문항 단위로 일반(2자몽) 또는 고급(8자몽) 모델 선택
  \item 학생 학습 계획표에 배정
\end{enumerate}

\begin{tipbox}{알아두면 좋은 점}
\begin{itemize}[leftmargin=1.2em]
  \item 본 기관 학습은 \textbf{본 기관 학생에게만} 노출 (다른 기관·본사 학생에게는 보이지 않음)
  \item 본사 PUBLIC 학습은 모든 기관 학생이 학습 가능
  \item AI 사용량은 \textbf{AI 자몽 지갑} 의 거래 이력에서 모두 확인 가능
  \item 학생이 직접 만든 OWN 학습은 본사가 별도 검수 (학생 매뉴얼 참고)
\end{itemize}
\end{tipbox}
